\textbf{Phần 1. Trắc nghiệm lựa chọn}
\setcounterpageref{ex}{0}

\begin{ex} %[4X9N52]
Dãy số liệu về chiều cao 4 bạn ($132\text{ cm}, 135\text{ cm}, 130\text{ cm}, 139\text{ cm}$) có giá trị lớn nhất là
\choice
{\True $139\text{ cm}$}
{$135\text{ cm}$}
{$132\text{ cm}$}
{$140\text{ cm}$}
\loigiai{
$139\text{ cm}$ lớn nhất.\\
\textbf{Đáp án của câu hỏi là A}
}
\end{ex}

\begin{ex} %[4X9N52]
Biểu đồ cột dùng để
\choice
{Vẽ tranh phong cảnh}
{Tính nhẩm phép nhân}
{Đo độ dài đoạn thẳng}
{\True Biểu diễn số liệu thống kê một cách trực quan}
\loigiai{
Biểu diễn số liệu thống kê trực quan.\\
\textbf{Đáp án của câu hỏi là D}
}
\end{ex}

\begin{ex} %[4X9H52]
Số trung bình cộng của dãy số liệu $12, 14, 16, 18$ là
\choice
{$14$}
{\True $15$}
{$16$}
{$17$}
\loigiai{
$(12 + 14 + 16 + 18) : 4 = 60 : 4 = 15$.\\
\textbf{Đáp án của câu hỏi là B}
}
\end{ex}

\begin{ex} %[4X9H52]
Tung đồng xu 20 lần có 12 lần ngửa. Số lần xuất hiện mặt sấp là
\choice
{$6$ lần}
{$10$ lần}
{$12$ lần}
{\True $8$ lần}
\loigiai{
$20 - 12 = 8$ lần.\\
\textbf{Đáp án của câu hỏi là D}
}
\end{ex}

\begin{ex} %[4X9H52]
Biểu đồ cột số điểm tốt của 3 tổ: Tổ 1 ($35$), Tổ 2 ($40$), Tổ 3 ($45$). Tổng số điểm tốt của cả 3 tổ là
\choice
{$110$}
{$130$}
{\True $120$}
{$115$}
\loigiai{
$35 + 40 + 45 = 120$.\\
\textbf{Đáp án của câu hỏi là C}
}
\end{ex}

\begin{ex} %[4X9H52]
Từ dữ liệu trên, trung bình mỗi tổ có bao nhiêu điểm tốt?
\choice
{$38$ điểm tốt}
{\True $40$ điểm tốt}
{$42$ điểm tốt}
{$45$ điểm tốt}
\loigiai{
$120 : 3 = 40$ điểm tốt.\\
\textbf{Đáp án của câu hỏi là B}
}
\end{ex}

\begin{ex} %[4X9V52]
Dãy số liệu số thóc 4 kho chứa ($25\text{ tấn}, 30\text{ tấn}, 35\text{ tấn}, 40\text{ tấn}$). Kho chứa nhiều nhất hơn kho chứa ít nhất bao nhiêu tấn thóc?
\choice
{\True $15\text{ tấn}$}
{$10\text{ tấn}$}
{$20\text{ tấn}$}
{$5\text{ tấn}$}
\loigiai{
$40 - 25 = 15$ tấn.\\
\textbf{Đáp án của câu hỏi là A}
}
\end{ex}

\begin{ex} %[4X9V52]
Trong một hộp có 4 bút bi xanh và 4 bút bi đỏ. Bạn Nam nhắm mắt lấy ra 1 chiếc bút. Khả năng lấy được bút bi xanh là
\choice
{Chắc chắn xảy ra}
{Không thể xảy ra}
{\True Bằng khả năng lấy được bút bi đỏ}
{Lớn hơn khả năng lấy bút bi đỏ}
\loigiai{
Số lượng 2 màu bằng nhau (4 chiếc) nên khả năng bằng nhau.\\
\textbf{Đáp án của câu hỏi là C}
}
\end{ex}

\begin{ex} %[4X9V52]
Gieo con xúc xắc 60 lần, mặt 6 chấm xuất hiện 12 lần. Số lần xuất hiện các mặt còn lại là
\choice
{$46$ lần}
{$50$ lần}
{$52$ lần}
{\True $48$ lần}
\loigiai{
$60 - 12 = 48$ lần.\\
\textbf{Đáp án của câu hỏi là D}
}
\end{ex}

\begin{ex} %[4X9V52]
Biểu đồ số học sinh 4 lớp: 4A ($35$), 4B ($38$), 4C ($34$), 4D ($37$). Tổng số học sinh của cả khối 4 là
\choice
{$140$ học sinh}
{$145$ học sinh}
{\True $144$ học sinh}
{$150$ học sinh}
\loigiai{
$35 + 38 + 34 + 37 = 144$ học sinh.\\
\textbf{Đáp án của câu hỏi là C}
}
\end{ex}

\begin{ex} %[4X9V52]
Trung bình mỗi lớp khối 4 có bao nhiêu học sinh?
\choice
{$35$ học sinh}
{\True $36$ học sinh}
{$37$ học sinh}
{$38$ học sinh}
\loigiai{
$144 : 4 = 36$ học sinh.\\
\textbf{Đáp án của câu hỏi là B}
}
\end{ex}

\begin{ex} %[4X9C52]
Một hộp có 10 quả bóng bàn gồm 3 màu: trắng, vàng, cam. Biết số bóng trắng gấp đôi số bóng vàng và số bóng cam là 1 quả. Bạn lấy ngẫu nhiên 1 quả bóng. Màu bóng có khả năng lấy được cao nhất là
\choice
{\True Bóng trắng (6 quả)}
{Bóng vàng}
{Bóng cam}
{Cả ba màu như nhau}
\loigiai{
Số bóng trắng $+$ vàng $= 10 - 1 = 9$ quả. Vàng: $9 : 3 = 3$ quả, Trắng: $3 \times 2 = 6$ quả.\\
Bóng trắng có 6 quả (nhiều nhất) nên khả năng lấy được cao nhất.\\
\textbf{Đáp án của câu hỏi là A}
}
\end{ex}

\textbf{Phần 2. Trắc nghiệm đúng sai}
\setcounterpageref{ex}{0}

\begin{ex} %[4X9H52]
Xét tính đúng sai về thống kê và xác suất:
\choiceTF
{\True Dãy số liệu thống kê giúp ta thu thập và xử lí thông tin số lượng.}
{\True Biểu đồ cột giúp so sánh các số liệu một cách trực quan.}
{Tung đồng xu 10 lần thì chắc chắn phải có 5 lần ngửa và 5 lần sấp.}
{\True Tung đồng xu 10 lần thì số lần ngửa và sấp có thể ngẫu nhiên khác nhau.}
\loigiai{
\begin{itemize}
    \item Dãy số liệu: Đúng.
    \item Biểu đồ cột trực quan: Đúng.
    \item Không chắc chắn đúng tỉ lệ 5-5 (xác suất là ngẫu nhiên). Sai.
    \item Số lần ngửa sấp là ngẫu nhiên: Đúng.
\end{itemize}
\begin{itemize}
    \item a) \textbf{Đúng}.
    \item b) \textbf{Đúng}.
    \item c) \textbf{Sai}.
    \item d) \textbf{Đúng}.
\end{itemize}
}
\end{ex}

\begin{ex} %[4X9V52]
Cho biểu đồ số ki-lô-gam giấy vụn thu gom của 3 lớp: 4A ($45\text{ kg}$), 4B ($50\text{ kg}$), 4C ($55\text{ kg}$). Xét tính đúng sai:
\choiceTF
{\True Lớp 4C thu gom được nhiều giấy vụn nhất.}
{\True Lớp 4A thu gom được ít giấy vụn nhất.}
{Cả 3 lớp thu gom được tất cả $160\text{ kg}$ giấy vụn.}
{\True Cả 3 lớp thu gom được tất cả $150\text{ kg}$ giấy vụn ($1{,}5$ tạ).}
\loigiai{
\begin{itemize}
    \item 4C: 55 kg (nhiều nhất). Đúng.
    \item 4A: 45 kg (ít nhất). Đúng.
    \item Tổng: $45 + 50 + 55 = 150\text{ kg} \neq 160\text{ kg}$. Ý 3 Sai, Ý 4 Đúng.
\end{itemize}
\begin{itemize}
    \item a) \textbf{Đúng}.
    \item b) \textbf{Đúng}.
    \item c) \textbf{Sai}.
    \item d) \textbf{Đúng}.
\end{itemize}
}
\end{ex}

\textbf{Phần 3. Trả lời ngắn}
\setcounterpageref{ex}{0}

\begin{ex} %[4X9H52]
Tìm trung bình cộng của 3 số liệu: $25, 35, 45$.
\shortans{35}
\loigiai{
\textbf{Đáp số:} $35$\\
$(25 + 35 + 45) : 3 = 35$.
}
\end{ex}

\begin{ex} %[4X9H52]
Tung đồng xu 30 lần có 16 lần ngửa. Số lần sấp là bao nhiêu?
\shortans{14}
\loigiai{
\textbf{Đáp số:} $14$\\
$30 - 16 = 14$ lần.
}
\end{ex}

\begin{ex} %[4X9V52]
Cho dãy số liệu điểm kiểm tra: $7, 8, 9, 8, 10, 8, 9$. Số điểm 8 xuất hiện bao nhiêu lần?
\shortans{3}
\loigiai{
\textbf{Đáp số:} $3$\\
3 lần.
}
\end{ex}

\begin{ex} %[4X9V52]
Biểu đồ số cây trồng của 3 tổ: $30, 40, 50$. Tính số cây trung bình mỗi tổ trồng được.
\shortans{40}
\loigiai{
\textbf{Đáp số:} $40$\\
$(30 + 40 + 50) : 3 = 40$ cây.
}
\end{ex}

\begin{ex} %[4X9V52]
Hiệu số cây giữa tổ trồng nhiều nhất (50) và tổ ít nhất (30) là bao nhiêu cây?
\shortans{20}
\loigiai{
\textbf{Đáp số:} $20$\\
$50 - 30 = 20$ cây.
}
\end{ex}

\begin{ex} %[4X9C52]
Một nhóm 4 học sinh có số đo chiều cao trung bình là $136\text{ cm}$. Khi có thêm bạn thứ 5 cao $141\text{ cm}$ tham gia nhóm thì chiều cao trung bình của cả 5 bạn là bao nhiêu xăng-ti-mét?
\shortans{137}
\loigiai{
\textbf{Đáp số:} $137$\\
Tổng 4 bạn: $136 \times 4 = 544\text{ cm}$.\\
Tổng 5 bạn: $544 + 141 = 685\text{ cm}$.\\
Trung bình 5 bạn: $685 : 5 = 137\text{ cm}$.
}
\end{ex}
